\chapter{Multiplayer}

\section{From One Agent to Many}

Every chapter of Part V so far has quietly favored the case of a single agent acting on the world. Chapter 15.3 and Chapter 27.6 each mentioned, in passing, that simultaneous proposals from different agents are resolved through joint admissibility rather than through a designated authority, and Chapter 28.4 established that human and AI agents enter the interaction layer through the same interface. None of these chapters actually worked through what happens once many agents, human and artificial, are acting on the same persistent field at overlapping moments, which is the ordinary condition of any world meant to be shared rather than merely visited. This chapter completes that picture.

\section{Commutative Proposals: The Easy Case}

It is worth starting with the case that requires no new machinery at all, because it is also the most common one in practice. Two agents, proposing changes to disjoint regions of the field at overlapping moments, one placing a bench in the courtyard, another inspecting a crack on the far side of the wall, are proposing commutative updates in Chapter 15's exact sense: each proposal projects and commits independently, in either order, with no need for the two agents to know about one another at all. Most of what a genuinely populated persistent world requires, at any given moment, falls into this category, and it is worth being honest that this case was already fully solved by Chapter 15. Multiplayer, as an engineering problem, is not primarily about this case. It is about the much rarer, much harder case the rest of this chapter addresses.

\section{Genuine Conflicts: When Proposals Collide}

Two agents proposing incompatible changes to the same region at overlapping moments, one moving a bench to the courtyard's east wall, another moving it, at nearly the same instant, to the west wall, cannot both commit. Chapter 15.3 stated the principle governing this case without working through its mechanics: resolution proceeds via joint admissibility rather than an arbitrary tie-break or a first-come-first-served rule. Concretely, this means the projection stage does not process the two proposals independently and then discover a conflict after the fact. It processes them jointly, computing the nearest point in C to the combined pair, and the result depends on what C actually permits: perhaps only one proposal's bench-placement survives intact while the other is refused outright, perhaps some partial reconciliation exists, a location between the two that satisfies neither agent's exact intent but violates nothing either proposed, if C happens to admit such a compromise. What matters is that the resolution is never arbitrary in the sense of depending on which proposal happened to arrive microseconds earlier. It depends only on what the field's own constraints, established since Chapter 18, actually permit jointly, and an agent whose proposal is refused under this scheme has not lost a race; it has proposed something that, given what another agent simultaneously and equally legitimately proposed, could not be jointly admitted.

\section{Beyond Ordinary Conflict: Adversarial and Pathological Patterns}

Section 29.3's machinery handles conflicts between individually legitimate proposals. It has nothing to say about a different and genuinely harder problem: a sequence of proposals, each individually admissible under \(\Pi_C\), that together constitute a pattern no shared world should tolerate. An agent that repeatedly proposes the same trivial, technically-admissible action at high frequency is not violating any single admissibility check, and yet is degrading the shared world for every other participant in it, consuming the write cycle's attention without contributing anything \(\Pi_C\) is equipped to object to. An agent whose proposals, while each individually harmless, form a pattern that another participant experiences as harassment is a harder case still, one where no single proposal is the problem and the problem is visible only in aggregate.

This is worth naming honestly as a genuinely different kind of check than anything Part IV or the earlier chapters of Part V developed. \(\Pi_C\) asks whether a proposal is consistent with the field's own constraints. It was never designed to ask whether a pattern of proposals, each consistent with those constraints taken alone, is healthy for the community of agents sharing the field. A separate monitoring layer is needed, one that watches patterns over time rather than checking single proposals in isolation, and the clearest existing precedent for this kind of layer, worth crediting directly, is the Custodian mechanism described in Christopher and Alexandra Chenoweth's MEM|8 architecture: a guard system that tracks repetition scores and psychological-safety signals across a stream of memory events, throttling or intervening when patterns cross thresholds that no single event would trigger on its own, including explicit mechanisms for preventing what that work calls repetition poisoning, an agent trapped in or exploiting a repetitive loop.

\section{A Second, Social Layer of Admissibility}

It is worth being precise about how this second layer relates to everything Part IV already built, because the two are complementary rather than competing. \(\Pi_C\), and the constraint manifold C behind it, governs field admissibility: whether a proposal is consistent with the persistent world's own accumulated structure. A Custodian-style monitor governs something this book will call social admissibility: whether a pattern of proposals, evaluated across a window of time rather than at a single instant, is consistent with a shared space multiple agents can continue to participate in. These are different questions, and a proposal can pass one while straining the other; the repeated trivial action from section 29.4 passes every field-level check and still degrades the shared world. This book will not attempt a full formal treatment of social admissibility comparable to Part IV's treatment of C; the concept is borrowed, deliberately, from an architecture built to address exactly this problem, and this chapter's contribution is locating where such a layer belongs in the write cycle this book has developed, downstream of \(\Pi_C\), monitoring the stream of what actually commits, rather than folding social concerns into C itself and overloading a mechanism that was built to answer a different question.

\section{Synchronization Without Consensus}

It is worth closing with a point this book's supporting technical material makes explicitly and that deserves to be stated plainly: none of the machinery this chapter has developed requires an explicit consensus protocol, a designated coordinator deciding whose proposal is authoritative. Different agents, particularly across a network with real latency, may hold temporarily divergent local views of the field, and this is not, on its own, a violation of anything this book has argued for. What is required is that these local views reconcile once commits actually propagate, in the same sense Chapter 21's gluing condition requires locally admissible patches to assemble into one coherent global section rather than requiring every observer to have identical information at every instant. Admissibility itself, checked at commit time, does the arbitration work a consensus protocol would otherwise be needed for; two agents' temporarily divergent views are reconciled not by a vote or a designated authority but by the same joint-admissibility computation section 29.3 already described. This chapter will not develop the full distributed-systems treatment such reconciliation requires at genuine scale, across many machines and significant network latency; that belongs to Chapter 30.

\section{Looking Ahead}

This chapter has treated multiplayer as a logical and architectural question: how conflicting proposals are resolved, how patterns no single proposal reveals are monitored, and why no explicit consensus protocol is required for either. It has not addressed the physical question of actually running this architecture across many machines, with real network latency, real hardware failure, and a world database too large for any single machine to hold. Chapter 30 takes up scaling directly, extending everything this Part has built to genuine distributed operation.
